\part {Using FishNET} {}

\chap Introductory Experiments

The back-propagation algorithm described in an earlier section was
applied to several problems using the software package fishNET.
Various input data was used, and the number of middle (intermediate)
layer \neus\ was varied.
The learning parameters $\alpha$ and $\varepsilon$ were also
modified.

Due to most of the work being carried out on a PC, there was an
upper limit to the number of \neus\ used in simulations.
The case dealt with in detail was extensively tested with parameter
variations.
The other examples here will be examined briefly to show how
versatile a \nn\ is.

The intent of this chapter is to illustrate
the initial experiments that were performed with fishNET to
prove that it worked as desired.
It is possible to use fishNET to do many different tasks simply by creating
the appropriate input and output data.
The software will then generate a generalised network, with no other
human input.
It is a remarkable tool.

Later chapters investigate further the effects of parameters on
network behaviour, and the effect of casualties within the network.

\sec The First Test --- O and X Recogniser

The first piece of test data that was created and run on the network
was a simple O and X recogniser.
There were~$9$ \neus\ (bits) on the input, $6 \to 9$ \neus\ in the middle
layer, and~$2$ output \neus\ (bits), each one of which signalled either a O
or a X.
The network converged to a solution in about~$100$ learning
iterations for~$7$ intermediate layer \neus; the speed of learning
versus number of intermediate \neus\ will be discussed later.

The input and output data is shown below
\def\tablerule{\noalign{\hrule}}
$$
\eqalign{\vbox{\tabskip=0pt \offinterlineskip
\halign{\strut#& \vrule#\tabskip=1em plus2em& \hfil#& \vrule#&
 \hfil#\hfil& \vrule#& \hfil#& \vrule#\tabskip=0pt\cr\noalign{\hrule}
&&\multispan5\hfil Input 1\hfil&\cr\noalign{\hrule height1.5pt}
&&1&&0&&1&\cr\noalign{\hrule}
&&0&&1&&0&\cr\noalign{\hrule}
&&1&&0&&1&\cr\noalign{\hrule}
}}}
\quad\eqalign{\vbox{\tabskip=0pt \offinterlineskip
\halign{\strut#& \vrule#\tabskip=1em plus2em& \hfil#& \vrule#&
 \hfil#\hfil& \vrule#& \hfil#& \vrule#\tabskip=0pt\cr\noalign{\hrule}
&&\multispan5\hfil Input 2\hfil&\cr\noalign{\hrule height1.5pt}
&&1&&1&&1&\cr\noalign{\hrule}
&&1&&0&&1&\cr\noalign{\hrule}
&&1&&1&&1&\cr\noalign{\hrule}
}}}
\qquad\eqalign{\vbox{\tabskip=0pt \offinterlineskip
\halign{\strut#& \vrule#\tabskip=1em plus2em& \hfil#\hfil& \vrule#&
 \hfil#\hfil& \vrule#\tabskip=0pt\cr\noalign{\hrule}
&&\multispan3\hfil Output 1\hfil&\cr\noalign{\hrule height1.5pt}
&&0&&1&\cr\noalign{\hrule}
}}}
\quad\eqalign{\vbox{\tabskip=0pt \offinterlineskip
\halign{\strut#& \vrule#\tabskip=1em plus2em& \hfil#\hfil& \vrule#&
 \hfil#\hfil& \vrule#\tabskip=0pt\cr\noalign{\hrule}
&&\multispan3\hfil Output 2\hfil&\cr\noalign{\hrule height1.5pt}
&&1&&0&\cr\noalign{\hrule}
}}}
$$
\centerline{Figure 8. Two simple input/output case pairs.}

To test if the network generalised, even for such a simple
case, input data not used for training was input (such as~$0.5$s
being used to replace some of the~$1$s in the shapes).
The network produced the correct results.

\sec A More Ambitious Test --- Shadow Encoding

Burr [25] discusses the use of a~$13$ segment shadow encoder as a
means of describing hand written data for input to a \nn.
He aimed to be able to recognise the alphabet as written by a
particular writer.
He took~$8$ samples of the alphabet written over a period of time,
and then trained the network on~4 of the samples, using the other~4
as test data to see if the network could successfully recognise
them.
His approach was remarkably fruitful, with a recognition rate of up
to~99\%.
The example given here was not that large.

Shadow encoding can be illustrated by figure 9,
from Burr [25].
\figure {9} {5} {Shadow encoding of the letter `S'}

Several attempts were made at using this shadow encoded alphabet
with fishNET.
Initially, 5 letters were encoded and the network was trained to
recognise them within a hundred or so learning iterations.
However, when the~5 letters were replaced by the full alphabet, the
computational time per learning iteration of the network rose by a factor
of more than~7 (since~5 times as many input/output cases, plus more
neurons in the middle layer).
Hence this network was never fully allowed to converge to a
solution.
Possible methods of speeding convergence are dealt with in a later
chapter.

\sec Dot Matrix Letter Encoding

This was the other major initial test experiment.
Shadow encoding, although partly removing some of the inherent problems of
\nn\ character recognition such as rotation, can't be easily
visualised.
The aim of this thesis is to provide qualitative judgements on
network behaviour, so for this reason a more easily visualised
encoding technique was sought.
Dot matrix encoding of several letters of the alphabet was chosen.

Various sizes (geometries) of the input layer were tried, and to make symmetric
letters easier to encode an odd number of \neus\ was always used in
both dimensions.
The number of \neus\ in the bottom layer was increased as the
processing speed increased due to the addition of the floating
point processor (8087).
First attempts were for a~$5\times 7$ grid of input \neus, but this
was too small for reasonable test data to be generated.
$9\times 11$ was also tried, before finally selecting~$13\times 15$.
This rather large number of \neus, (195 in the bottom layer,) gives
a dot matrix encoded letter that is of reasonable quality, while
still being easily computable in a network on a PC with
an~8087.

This example is extensively discussed in the next chapter, and is
used throughout the thesis for experimental purposes.

